La matriz de escalamiento es utilizada para definir cómo se deben comunicar los involucrados del proyecto según el tipo de comunicación que se requiera transmitir y el rol de los mismos. Este proceso asegura que los incidentes que no puedan resolverse a un nivel inferior sean gestionados de manera oportuna por la cadena de mando adecuada.

\begin{table}[H]
\centering
\scriptsize 
\setlength{\tabcolsep}{5pt}
\renewcommand{\arraystretch}{1.8}
\begin{tabular}{|p{3cm}|p{3.5cm}|p{4cm}|p{3.5cm}|}
\hline
\rowcolor[HTML]{003366} 
\textcolor{white}{\textbf{NIVEL DE ESCALAMIENTO}} & 
\textcolor{white}{\textbf{TIPO DE COMUNICACIÓN}} & 
\textcolor{white}{\textbf{NOMBRE}} & 
\textcolor{white}{\textbf{ROL}} \\ \hline

\textbf{Nivel 1: Técnico} & Comunicación Interactiva (Mensajería instantánea, discusiones ad hoc) & [Jhoan Esteban Corrales] & Líder de Desarrollo / Coordinador Técnico \\ \hline

\textbf{Nivel 2: Operativo} & Comunicación Formal (Informes, correos electrónicos de seguimiento)  & [Juan Camilo Velasquez Tobo] & Analista Funcional / Especialista de Seguridad \\ \hline

\textbf{Nivel 3: Gestión} & Gestión de expectativas y resolución de conflictos en reuniones & [Jorge Luis Araque] & Project Manager (Director del Proyecto) \\ \hline

\textbf{Nivel 4: Estratégico} & Comunicación Oficial (Actas de comité, informes anuales) & [Julio César Lopez] & Patrocinador / Comité Directivo \\ \hline

\textbf{Nivel 5: Externo} & Comunicación Escrita Oficial (Oficios institucionales)  & [Representante RUNT/MinTransporte] & Autoridad Regulatoria / Stakeholder Externo \\ \hline

\end{tabular}
\end{table}

\subsubsection*{Consideraciones del Procedimiento de Escalamiento}
\begin{itemize}
    \item \textbf{Criterio de Urgencia:} La selección de la tecnología y el nivel de escalamiento dependerá de la urgencia de la necesidad de información.
    \item \textbf{Responsabilidad del Emisor:} El emisor en cada nivel es responsable de asegurar que la información comunicada sea clara, completa y comprendida correctamente antes de escalar al siguiente nivel.
    \item \textbf{Retroalimentación:} Una vez descodificado el mensaje de escalamiento, el receptor debe proporcionar una respuesta o retroalimentación clara sobre la resolución del incidente.
    \item \textbf{Registro de Incidentes:} Todos los eventos escalados deben ser documentados en el registro de incidentes como entrada para el proceso de control de comunicacione.
\end{itemize}